\input{../tex-inputs/latex-leseansicht-vorspann}

\section[Emil Marriot, Rudolf Kobatsch und Anton Peter {[}?{]} Mataja an Arthur Schnitzler, 20. 7. 1902]{L04337 Emil Marriot, Rudolf Kobatsch und Anton Peter [?] Mataja an Arthur
               Schnitzler, 20. 7. 1902}
\nopagebreak\mylabel{L04337v}
\rehead{ }\normalsize\beginnumbering\briefempfaengerindex{Schnitzler, Arthur@\textsc{Schnitzler, Arthur}!zzzMataja, Anton Peter@\emph{von Anton Peter Mataja}!1902-07-201@{20. 7. 1902}|(be}\briefempfaengerindex{Schnitzler, Arthur@\textsc{Schnitzler, Arthur}!zzzKobatsch, Rudolf@\emph{von Rudolf Kobatsch}!1902-07-201@{20. 7. 1902}|(be}\briefempfaengerindex{Schnitzler, Arthur@\textsc{Schnitzler, Arthur}!zzzMarriot, Emil@\emph{von Emil Marriot}!1902-07-201@{20. 7. 1902}|(be}
\toendnotes[C]{\smallbreak\pagebreak[2]}
\correspDesc{Versand  durch Emil Marriot, Rudolf Kobatsch, Anton Peter Mataja am 20. 7. 1902 in Mödling
\newline{}Übermittlung  am 20. 7. 1902 in Mödling
\newline{}Zustellung  am 21. 7. 1902 in Wien
\newline{}Erhalt  durch Arthur Schnitzler am 21. 7. 1902 in Wien}\toendnotes[C]{\smallbreak}
\Standort{DLA, A:Schnitzler, HS.1985.1.4008.}
\physDesc{Bildpostkarte, 255 Zeichen
\newline{}Handschrift: Bleistift, lateinische Kurrent
\newline{}Versand: 1) Stempel: »\nobreak{}\oindex{Mödling@\textbf{Mödling}, \emph{Hauptstadt}|pwk}Mödling, 20. {[}7. 02{]}, 8–9\nobreak{}«.   2) Stempel: »\nobreak{}\oindex{IX., Alsergrund@\textbf{IX., Alsergrund}, \emph{Verwaltungsgebiet}|pwk}9/3 Wien 72, 21. 7. 02, 8.V, Bestellt\nobreak{}«. }\toendnotes[C]{\smallbreak}\pstart{}{\pb}Herrn\pend{}\pstart{}D\textsuperscript{r.} Arthur Schnitzler\pend{}\pstart{}Wien IX\oindex{IX., Alsergrund@\textbf{IX., Alsergrund}, \emph{Verwaltungsgebiet}|pw}\pend{}\pstart{}Frankgasse 1\oindex{Wien@\textbf{Wien}!IX., Alsergrund@\textbf{IX., Alsergrund}!Frankgasse 1@\textbf{Frankgasse 1}, \emph{Wohngebäude}|pw}.\pend{}{\bigskip}
\pstart
           \noindent{}\centering{}{\pb}\textcolor{gray}{\textbf{Gruss aus Vorderbrühl,
                     Meierei\oindex{Liechtensteinische Meierei@\textbf{Liechtensteinische Meierei}, \emph{Gastgewerbegebäude}|pw}}}\pend
           \vspace{1em}
\pstart
           \noindent{}{\pb}Ich bin hier mit Herrn Karl Kraus\pwindex{Kraus, Karl 28.\,4.\,1874 Jičín – 12.\,6.\,1936 Wien@\textsc{Kraus, Karl} (28.\,4.\,1874 Jičín – 12.\,6.\,1936 Wien), \emph{Schriftsteller, Publizist, Schriftsteller}|pw} und freue mich, dass \label{K_L04337-1v}\edtext{heute hier das »Abschiedssouper\pwindex{Schnitzler, Arthur 15. 5. 1862 Wien – 21. 10. 1931 ebd.@\textsc{Schnitzler, Arthur} (15. 5. 1862 Wien – 21. 10. 1931 ebd.), \emph{Schriftsteller, Mediziner}!Abschiedssouper@\strich\emph{Abschiedssouper}|pw}« gegeben}{\lemma{\textnormal{\emph{heute … gegeben}}}\Cendnote{\textnormal{Die
                     Aufführung ist derzeit nicht belegt. Am 18. 7. 1902 führte das \emph{Kurtheater in Kaltenleutgeben}\orgindex{Kurtheater in Kaltenleutgeben@Kurtheater in Kaltenleutgeben|pwk} im 10 Kilometer entfernten Kaltenleutgeben\oindex{Kaltenleutgeben@\textbf{Kaltenleutgeben}, \emph{Hauptstadt}|pwk}{ }\emph{Abschiedssouper}\pwindex{Schnitzler, Arthur 15. 5. 1862 Wien – 21. 10. 1931 ebd.@\textsc{Schnitzler, Arthur} (15. 5. 1862 Wien – 21. 10. 1931 ebd.), \emph{Schriftsteller, Mediziner}!Abschiedssouper@\strich\emph{Abschiedssouper}|pwk} 
                     gemeinsam mit \emph{Das Versprechen hinter’m Herd (Eine Szene aus den
                        österreichischen Alpen)}\pwindex{Baumann, Alexander 7.\,2.\,1814 Wien – 25.\,12.\,1857 Graz@\textsc{Baumann, Alexander} (7.\,2.\,1814 Wien – 25.\,12.\,1857 Graz)!Versprechen hinter’m Herd. Eine Szene aus den österreichischen Alpen@\strich\emph{Das Versprechen hinter’m Herd. Eine Szene aus den österreichischen Alpen}|pwk} und Chansons\eventindex{Kaltenleutgeben@\textbf{Kaltenleutgeben}!Aufführung von Abschiedssouper; Chansons; Das Versprechen hinter’m Herd, 18.7.1902@Aufführung von Abschiedssouper; Chansons; Das Versprechen hinter’m Herd, 18.7.1902|pwk} auf. Dieser Abend dürfte am 20. 7. 1902 
                     wiederholt worden sein, aber wo genau, ist unklar.}}}\label{K_L04337-1} wird – Wo sind Sie
                  \textcolor{gray}{ſelbs}t?\pend
           \pstart \spacefill\mbox{Marriot}\pend{}
\pstart
           20 Juni 1902\pend
           
\pstart
           Morgen reise ich nach Fieberbrunn
                     in Tirol\oindex{Fieberbrunn@\textbf{Fieberbrunn}, \emph{Verwaltungsgebiet}|pw}  ab.\pend
           \selectlanguage{ngerman}\vspace{1em}
\pstart
           \noindent{}{[}hs. Kobatsch:{]} Ein Verehrer\pend
           \pstart \spacefill\mbox{D\textsuperscript{r}Kobatsch}\pend{}\selectlanguage{ngerman}\vspace{1em}\pstart {[}hs. Mataja:{]} \spacefill\mbox{\label{K_L04337-2v}\edtext{Mataja}{\lemma{\textnormal{\emph{Mataja}}}\Cendnote{\textnormal{Die Zuordnung der Unterschrift an den Vater von Emil Marriot\pwindex{Marriot, Emil 20.\,11.\,1855 Wien – 5.\,5.\,1938 ebd.@\textsc{Marriot, Emil} (20.\,11.\,1855 Wien – 5.\,5.\,1938 ebd.), \emph{Schriftstellerin}|pwk} ist unsicher. 
                     Als Argument dafür lässt sich vor allem anführen, dass nur das Familien›oberhaupt‹ die Selbstverständlichkeit in Anspruch nehmen konnte, ohne Vorname zu unterschreiben.
                     Es könnte sich nichtsdestotrotz auch um ein anderes Familienmitglied handeln. }}}\label{K_L04337-2}}\pend{}\selectlanguage{ngerman}\endnumbering\briefempfaengerindex{Schnitzler, Arthur@\textsc{Schnitzler, Arthur}!zzzMataja, Anton Peter@\emph{von Anton Peter Mataja}!1902-07-201@{20. 7. 1902}|)be}\briefempfaengerindex{Schnitzler, Arthur@\textsc{Schnitzler, Arthur}!zzzKobatsch, Rudolf@\emph{von Rudolf Kobatsch}!1902-07-201@{20. 7. 1902}|)be}\briefempfaengerindex{Schnitzler, Arthur@\textsc{Schnitzler, Arthur}!zzzMarriot, Emil@\emph{von Emil Marriot}!1902-07-201@{20. 7. 1902}|)be}\mylabel{L04337h}
\begin{anhang}
\end{anhang}\newcommand{\dateiname}{L04337}\newcommand{\titel}{Emil Marriot, Rudolf Kobatsch und Anton Peter [?] Mataja an Arthur Schnitzler, 20. 7. 1902}\newcommand{\editorInnen}{Herausgegeben von Jahnke, SelmaMüller, Martin Anton}\input{../tex-inputs/latex-leseansicht-abspann}
